%=========== ΛΥΣΗ - 53ο ΘΕΜΑ Γ ===========
\begin{thema}{Γ}
\begin{erwthma}
\item Έστω $a$ η πλευρά του τετραγώνου και $b$ η πλευρά του ισόπλευρου τριγώνου.
Αφού $x$ είναι η περίμετρος του τετραγώνου, έχουμε
\[4a=x\Rightarrow a=\frac{x}{4}\]
Το υπόλοιπο σύρμα έχει μήκος $10-x$ και αποτελεί την περίμετρο του ισόπλευρου τριγώνου, άρα
\[3b=10-x\Rightarrow b=\frac{10-x}{3}\]
Επειδή κατασκευάζονται και τα δύο σχήματα, είναι
\[0<x<10\]

\item Το εμβαδόν του τετραγώνου είναι
\[E_1=a^2=\left(\frac{x}{4}\right)^2=\frac{x^2}{16}\]
Το εμβαδόν του ισόπλευρου τριγώνου είναι
\[E_2=\frac{\sqrt{3}}{4}b^2=\frac{\sqrt{3}}{4}\left(\frac{10-x}{3}\right)^2=\frac{\sqrt{3}(10-x)^2}{36}\]
Άρα το συνολικό εμβαδόν των δύο σχημάτων είναι
\[E(x)=E_1+E_2=\frac{x^2}{16}+\frac{\sqrt{3}(10-x)^2}{36}\ ,\ 0<x<10\]

\item Η $E$ είναι παραγωγίσιμη στο $(0,10)$ με
\[E'(x)=\frac{x}{8}-\frac{\sqrt{3}(10-x)}{18}\]
Βρίσκουμε τη ρίζα της $E'$:
\begin{align*}
E'(x)=0
&\Leftrightarrow \frac{x}{8}=\frac{\sqrt{3}(10-x)}{18}\\
&\Leftrightarrow 18x=8\sqrt{3}(10-x)\\
&\Leftrightarrow (18+8\sqrt{3})x=80\sqrt{3}\\
&\Leftrightarrow x=\frac{40\sqrt{3}}{9+4\sqrt{3}}
\end{align*}
Επιπλέον
\[E''(x)=\frac{1}{8}+\frac{\sqrt{3}}{18}>0\]
οπότε η $E$ είναι κυρτή στο $(0,10)$. Άρα το ελάχιστο συνολικό εμβαδόν επιτυγχάνεται για
\[x_0=\frac{40\sqrt{3}}{9+4\sqrt{3}}\]

Για τον λόγο της πλευράς του τετραγώνου προς την πλευρά του τριγώνου, από τη σχέση $E'(x_0)=0$ έχουμε
\[\frac{x_0}{8}=\frac{\sqrt{3}(10-x_0)}{18}\Rightarrow \frac{x_0}{10-x_0}=\frac{4\sqrt{3}}{9}\]
Άρα
\[\frac{a}{b}=\frac{\frac{x_0}{4}}{\frac{10-x_0}{3}}=\frac{3x_0}{4(10-x_0)}=\frac{3}{4}\cdot\frac{4\sqrt{3}}{9}=\frac{\sqrt{3}}{3}\]

\item Εφόσον ζητείται να κοπεί το σύρμα σε δύο τμήματα, πρέπει
\[0<x<10\]
Η $E$ είναι κυρτή στο $(0,10)$, άρα δεν παρουσιάζει μέγιστο σε εσωτερικό σημείο.
Για να δούμε προς ποιο άκρο μεγαλώνει περισσότερο το εμβαδόν, εξετάζουμε τη συνεχή επέκταση της $E$ στο $[0,10]$:
\[E(0)=\frac{\sqrt{3}\cdot 10^2}{36}=\frac{25\sqrt{3}}{9}\]
και
\[E(10)=\frac{10^2}{16}=\frac{25}{4}\]
Επειδή
\[\frac{25}{4}>\frac{25\sqrt{3}}{9}\]
το μεγαλύτερο δυνατό άνω φράγμα του εμβαδού είναι
\[\frac{25}{4}\]
και προσεγγίζεται όταν $x\to 10^-$, δηλαδή όταν δίνουμε σχεδόν όλο το σύρμα στο τετράγωνο και ένα πολύ μικρό τμήμα στο τρίγωνο.

Άρα, αν απαιτείται να κατασκευαστούν και τα δύο σχήματα, μέγιστο εμβαδόν δεν επιτυγχάνεται.
Αν επιτρέπεται εκφυλισμένη περίπτωση, τότε παίρνουμε $x=10$, δηλαδή χρησιμοποιούμε όλο το σύρμα για το τετράγωνο.
\end{erwthma}
\end{thema}
%=========== ΤΕΛΟΣ ΛΥΣΗΣ - 53ο ΘΕΜΑ Γ ===========
